\documentclass{article}
%\usepackage{esint}
\usepackage{booktabs}
\usepackage[portuguese]{babel}
\usepackage{macros}

\begin{document}

\tableofcontents

\section{Área 1}

\subsection{Lei de Coulomb}
\subsection{Campo Elétrico}
\subsection{Lei de Gauss}

\section{Área 2}

\subsection{Potencial Elétrico}
\subsection{Capacitância}
\subsection{Circuitos DC}


\section{Área 3}

\subsection{Campo Magnético}

A Lei de Biot-Savart descreve o campo magnético \(\mathbf{B}\) gerado por uma corrente elétrica \(I\) em um ponto \(P\) a uma distância \(\mathbf{r}\) do elemento de corrente \(d\mathbf{l}\):

\begin{equation}\label{eq:lei_biot_savart}
    \boxed{\mathbf{B}(r) = \frac{\mu_0}{4\pi} \int \frac{I d\mathbf{l} \times \mathbf{r}}{r^3}}
\end{equation} onde \(\mu_0\) é a permeabilidade magnética do vácuo, \(d\mathbf{l}\) é o elemento de corrente e \(\mathbf{r}\) é o vetor que aponta do elemento de corrente ao ponto \(P\).

Podemos utilizar a Lei de Biot-Savart para calcular o campo magnético gerado por diversas configurações de corrente elétrica, como espiras circulares e fios retilíneos infinitos.

\begin{itemize}
    \item [(a)] \textbf{Campo magnético de uma espira circular}:
        \begin{equation}
            \boxed{\mathbf{B}(z) = \frac{\mu_0 I R^2}{2(R^2 + z^2)^{3/2}} \hat{\mathbf{z}}}
        \end{equation} onde \(R\) é o raio da espira, \(I\) é a corrente elétrica e \(z\) é a distância do ponto ao plano da espira.
    \item [(b)] \textbf{Campo magnético ao redor de um fio retilíneo infinito}:
        \begin{equation}
            \boxed{\mathbf{B}(r) = \frac{\mu_0 I}{2\pi r}\hat{\boldsymbol{\varphi}}}
        \end{equation} onde \(r\) é a distância do ponto ao fio.
\end{itemize}

\subsection{Lei de Ampère}

A Lei de Ampère é uma das equações de Maxwell e relaciona o campo magnético \(\mathbf{B}\) com a corrente elétrica \(I\) que atravessa uma curva fechada \(C\), comumente chamada de curva amperiana. A lei é expressa pela seguinte equação:

\begin{equation}
    \boxed{\oint_{C} \mathbf{B} \cdot d\mathbf{l} = \mu_0 I_{C}}
\end{equation} onde \(I_{C}\) é a corrente elétrica que atravessa a superfície delimitada pela curva amperiana \(C\).

Em problemas envolvendo simetria, podemos aplicar a Lei de Ampère para calcular o campo magnético gerado por diferentes condutores, como fios retilíneos infinitos, toróides e solenoides.

\begin{itemize}
    \item [(a)] \textbf{Campo magnético de um fio retilíneo infinito}:\\
        \begin{align}
                \boxed{\mathbf{B}(r) = \begin{cases}
                    \frac{\mu_0 I}{2\pi} \frac{r}{R^2} \hat{\boldsymbol{\varphi}} \quad &(0 \leq r \leq R) \\
                    \frac{\mu_0 I}{2\pi} \frac{1}{r} \hat{\boldsymbol{\varphi}} \quad &(R \leq r)
                \end{cases}}
        \end{align} onde \(R\) é o raio do condutor e \(r\) é a distância do ponto ao fio.

    \item [(b)] \textbf{Campo magnético de um toróide}:
        \begin{align}
            \boxed{\mathbf{B}(r) = \begin{cases}
            		\frac{\mu_0 N I}{2\pi r}\hat{\boldsymbol{\theta}} \quad &(a < r < b) \\
            		0 \quad &(r < a \,\text{ ou }\, r > b)
            	\end{cases}}
        \end{align} onde \(N\) é o número de espiras do toróide de raios interno \(a\) e externo \(b\) e \(r\) é a distância do ponto ao eixo do toróide.

    \item [(c)] \textbf{Campo magnético de um solenoide}:
        \begin{align}
            \boxed{\mathbf{B} = \begin{cases}
            		\mu_0 n I\hat{\mathbf{z}} \quad &(\text{dentro do solenoide}) \\
            		0 \quad &(\text{fora do solenoide})
            	\end{cases}}
        \end{align} onde \(n\) é o número de espiras por unidade de comprimento do solenoide.
\end{itemize}

A \emph{força por unidade de comprimento} entre dois condutores paralelos infinitos é dada por:

\begin{equation}
    \boxed{f = I_1 B_2 = \frac{\mu_0 I_1 I_2}{2\pi d}}
\end{equation} onde \(I_1\) e \(I_2\) são as correntes nos condutores e \(d\) é a distância entre eles.

\section{Área 4}

\subsection{Indução Eletromagnética}

A Lei da Indução de Faraday afirma que a força eletromotriz induzida em um circuito fechado é proporcional à taxa de variação do fluxo magnético através do circuito.

\begin{equation}\label{eq:lei_faraday}
    \boxed{\varepsilon = -\frac{d\Phi_{B}}{dt}}
\end{equation} onde \(\varepsilon\) é a força eletromotriz induzida e \(\Phi_{B}\) é o fluxo magnético que atravessa o circuito. A unidade de fluxo magnético é o weber (Wb), que é equivalente a \(1 \, \text{Wb} \equiv 1 \, \text{T}\cdot\text{m}^2\).

O sinal negativo de \eqref{eq:lei_faraday} indica que a corrente induzida tem um sentido tal que o campo magnético gerado por ela se opõe à variação do fluxo magnético que a causou, conforme enunciado pela Lei de Lenz.

\subsection{Indutância}

Indutor é um componente elétrico que armazena energia na forma de um campo magnético quando uma corrente elétrica passa por ele. A indutância \(L\) de um indutor é uma medida de sua capacidade de induzir uma força eletromotriz (fem) em resposta a uma variação da corrente elétrica. A indutância é definida como

\begin{equation}
    \boxed{L \equiv \frac{N\Phi_{B}}{I}}
\end{equation} onde \(N\) é o número de espiras, \(\Phi_{B}\) é o fluxo magnético que atravessa cada espira e \(I\) é a corrente elétrica. A unidade de indutância é o henry (H), que é equivalente a \(1 \, \text{H} \equiv 1 \, \text{Wb}/\text{A}\).

Dada a força eletromotriz induzida num circuito \(\displaystyle \varepsilon = - \frac{d\Phi}{dt}\), a potência que deve ser fornecida para manter a corrente \(i\) é

\begin{equation}
    \frac{dW}{dt} = - \varepsilon i = + \frac{d\Phi}{dt} i = + Li \frac{di}{dt}
\end{equation} onde \(L\) é a autoindutância do circuito. Ignorando a perda por efeito Joule, a energia armazenada no circuito atravessado por uma corrente \(I\) é

\begin{equation}
    \boxed{U = \frac{1}{2} L I^2}
\end{equation}

A energia magnética armazenada por um solenoide, quando atravessado por uma corrente \(I\), é

\begin{equation}
    U = \frac{1}{2} L I^2 = \frac{1}{2}\mu_0 \left(NI\right)^2 \frac{S}{l} = \frac{1}{2\mu_0} \left(\frac{NI}{l}\right)^2 Sl = \frac{1}{2\mu_0} B^2 V
\end{equation} onde \(n = N/l\) é o número de espiras por unidade de comprimento do solenoide, \(B\) é o campo dentro do solenoide e \(V = Sl\) é o volume interno a ele. Essa energia está armazenada no campo magnético, com \emph{densidade de energia magnética}

\begin{equation}
    \boxed{u_m = \frac{1}{2\mu_0}\mathbf{B}^2}
\end{equation} da mesma forma que a \emph{densidade de energia elétrica} é dada por \(\displaystyle u = \frac{\varepsilon_0}{2}\mathbf{E}^2\). Se tivermos simultaneamente um campo elétrico e um magnético, a densidade de energia eletromagnética total do campo é

\begin{equation}
    \boxed{u = u_e + u_m = \frac{\varepsilon_0}{2}\mathbf{E}^2 + \frac{1}{2\mu_0} \mathbf{B}^2}
\end{equation}

\subsection{Circuitos RC e RL}

\begin{itemize}
    \item [(a)] \textbf{Circuito RC}\\
        Dado um circuito com capacitor inicialmente descarregado e ligado a uma bateria de fem \(\varepsilon\), sendo \(R\) a resistência do circuito (incluindo a resistência interna da bateria), a corrente \(i(t)\) no circuito é dada por
            \begin{equation}
                \boxed{i(t) = \frac{\varepsilon}{R} \left(e^{-\frac{t}{\tau_C}}\right)}
            \end{equation} onde \(\tau_C \equiv RC\) é a constante de tempo do circuito. Vemos que a corrente decresce exponencialmente com o tempo, tendendo a zero ao \(t \to \infty\). Já a carga \(q(t)\) no capacitor é dada por

            \begin{equation}
                \boxed{q(t) = C\varepsilon \left(1 - e^{-\frac{t}{\tau_C}}\right)}
            \end{equation} onde se nota que para \(t >> \tau_C\), a carga final do capacitor é \(C\varepsilon\).

    \item [(b)] \textbf{Circuito RL}\\
        Para um circuito RL, onde \(L\) é a indutância do indutor, a corrente \(i(t)\) é dada por

        \begin{equation}
            \boxed{i(t) = \frac{\varepsilon}{R} \left(1 - e^{-\frac{t}{\tau_L}}\right)}
        \end{equation} onde \(\displaystyle \tau_L \equiv \frac{L}{R}\) é a constante de tempo do circuito RL.
    \end{itemize}

\subsection{Oscilações Eletromagnéticas}

\subsubsection{Oscilações Livres: Circuito LC}

Consideramos um circuito LC ideal, em que há um capacitor de capacitância \(C\) e um indutor de autoindutância \(L\) conectados. Como desprezamos a resistência, não há dissipação de energia, e a energia total do circuito é conservada. Pela 1ª Lei de Kirchhoff, obtemos

\begin{equation}
    \frac{d^2I}{dt^2} + \omega_0^2 I = 0
\end{equation} onde \(\displaystyle \omega_0 \equiv \frac{1}{\sqrt{LC}}\) é a frequência angular das oscilações livres do circuito LC. A solução geral dessa equação diferencial é

\begin{equation}\label{eq:solucao_oscilador_livre}
    \boxed{I(t) = A \cos(\omega_0 t + \varphi)}
\end{equation} onde \(A\) é a amplitude da corrente e \(\varphi\) é a fase inicial. Integrando \eqref{eq:solucao_oscilador_livre}, obtemos a carga \(Q(t)\) no capacitor:

\begin{equation}
    \boxed{Q(t) = \frac{A}{\omega_0} \sin(\omega_0 t + \varphi)}
\end{equation} com duas constantes arbitrárias

\begin{align}
    \begin{aligned}
        I(0) &= I_0 = A \cos \varphi \\
        Q(0) &= Q_0 = \frac{A}{\omega_0} \sin(\varphi)
        \end{aligned}
    \qquad
    \begin{aligned}\begin{cases}
        A &= \sqrt{I_0^2 + \omega_0^2 Q_0^2} \\
        \varphi &= \arctan \left( \frac{\omega_0 Q_0}{I_0} \right)
    \end{cases}\end{aligned}
\end{align} o que determinam \(A\) e \(\varphi\) em função das condições iniciais \(I(0)\) e \(Q(0)\). Para \(I_0 = 0\) e carga inicial contida no capacitor, temos \(A = \omega_0 Q_0\) e \(\varphi = \frac{\pi}{2}\).

A energia armazenada no capacitor (inteiramente confinada no campo elétrico entre as placas do capacitor) no instante \(t\) é (usando \(\omega_0^2 \equiv 1/LC\))

\begin{equation}
    \boxed{U_C(t) = \frac{Q^2(t)}{2C} = \frac{A^2}{2\omega_0^2 C} \sin^2(\omega_0 t + \varphi) = \frac{1}{2}LA^2 \sin^2(\omega_0 t + \varphi)}
\end{equation} e a energia armazenada no indutor no instante \(t\) é

\begin{equation}
    \boxed{U_L(t) = \frac{1}{2}LI^2(t) = \frac{1}{2}LA^2 \cos^2(\omega_0 t + \varphi)}
\end{equation} que leva a uma energia total 

\begin{equation}
    \boxed{U(t) = U_C(t) + U_L(t) = \frac{1}{2}LA^2 = \frac{1}{2}\frac{A^2}{\omega_0^2 C}}
\end{equation}

Como \(\sin(x + \frac{\pi}{2}) = \cos x\), vemos que a corrente \(I\) está \emph{adiantada} de \(\pi/2\), na fase, em relação à carga \(Q\) no capacitor (está \emph{em quadratura}). Tanto a corrente quanto a carga trocam de sinal (sentido) a cada hemiciclo. Portanto, a energia oscila entre energia elétrica e energia magnética, mantendo constante a \emph{energia eletromagnética total}.

\begin{table}[h]
    \centering
    \caption{Comparação entre oscilador mecânico e oscilador L-C}
    \label{tab:osciladores}
    \begin{tabular}{l l}
    \toprule
    \textbf{Oscilador mecânico} & \textbf{Oscilador L-C} \\
    \midrule
    $\displaystyle m \frac{d^2 x}{dt^2} + kx = 0$ & $\displaystyle L \frac{d^2 Q}{dt^2} + \frac{1}{C} Q = 0$ \\
    $\displaystyle x, \, m, \, k, \, \omega_0 = \sqrt{\frac{k}{m}}$ & $\displaystyle Q, \, L, \, \frac{1}{C}, \, \omega_0 = \frac{1}{\sqrt{LC}}$ \\
    $\displaystyle v = \frac{dx}{dt}$ & $\displaystyle I = \frac{dQ}{dt}$ \\
    Energia cinética: $\displaystyle T = \frac{1}{2} mv^2$ & Energia magnética: $\displaystyle U_M = \frac{1}{2} LI^2$ \\
    Energia potencial: $\displaystyle V = \frac{1}{2} kx^2$ & Energia elétrica: $\displaystyle U_E = \frac{Q^2}{2C}$ \\
    \bottomrule
    \end{tabular}
\end{table}

\subsubsection{Oscilações Amortecidas: Circuito RLC}

Agora numa situação mais realista, consideramos um circuito RLC, onde \(R\) é a resistência do circuito. A equação diferencial que descreve o circuito é

\begin{equation}
    \frac{d^2Q}{dt^2} + \frac{R}{L} \frac{dQ}{dt} + \frac{1}{LC} Q = 0
\end{equation} substituindo \(\omega_0 \equiv \sqrt{1/LC}\) e definindo \(\gamma \equiv R/2L\), obtemos

\begin{equation}
    \frac{d^2Q}{dt^2} + 2\gamma \frac{dQ}{dt} + \omega_0^2 Q = 0
\end{equation} cuja solução geral é

\begin{equation}\label{eq:solucao_oscilador_amortecido}
    \boxed{Q(t) = A e^{-\gamma t} \cos(\omega_1 t + \varphi)}
\end{equation} onde \(\displaystyle \omega_1 \equiv \sqrt{\omega_0^2 - \gamma^2}\in\mathbb{R}\) é a frequência angular das oscilações amortecidas. A constante de amortecimento \(\gamma\) (atrito), introduzida pela resistência do circuito, determina a taxa de decaimento da amplitude da carga.

\subsubsection{Oscilações Forçadas: Circuito RLC com Fonte Alternada}

Quando uma fonte de tensão alternada é conectada a um circuito RLC, temos oscilações forçadas. A fonte fornece energia continuamente, mantendo as oscilações. Considerando uma fonte senoidal

\begin{equation}
    \varepsilon(t) = \varepsilon_0 \sin(\omega t)
\end{equation} a equação diferencial que descreve o circuito RLC forçado é

\begin{equation}
    \frac{d^2Q}{dt^2} + \frac{R}{L} \frac{dQ}{dt} + \frac{1}{LC} Q = \varepsilon_0 \sin(\omega t)
\end{equation} cuja solução geral possui duas partes: a solução transiente, que decai com o tempo, e a solução estacionária (ou permanente), que oscila na mesma frequência da fonte. A corrente no regime permanente tem a forma

\begin{equation}
    \boxed{I(t) = I_0 \sin(\omega t - \varphi)}
\end{equation} onde \(I_0\) é a amplitude da corrente e \(\varphi\) é a fase de defasagem entre a corrente e a tensão da fonte:

\begin{align}
    &I_0 = \frac{\varepsilon_0}{\sqrt{R^2 + \left(\frac{1}{\omega L} - \omega C\right)^2}} \\
    &\tan(\varphi) = \frac{\omega L - \frac{1}{\omega C}}{R}
\end{align} quando \(\displaystyle \omega = \omega_0 \equiv \frac{1}{\sqrt{LC}}\), temos ressonância, onde a amplitude da corrente é máxima.

\subsection{Circuitos AC}

\subsubsection{Circuito Resistivo Puro}

Consideramos um circuito resistivo puro, com uma fonte de tensão alternada \(V(t) = V_0 \sin(\omega t)\). Aplicando a 1ª Lei de Kirchhoff, temos

\begin{equation}
    - V_0 \sin(\omega t) + R i(t) = 0
\end{equation} isolando \(i(t)\), obtemos

\begin{equation}
    \boxed{i(t) = \frac{V_0}{R} \sin(\omega t) = I_0 \sin(\omega t)}
\end{equation} onde \(\displaystyle I_0 = V_0/R\) é a amplitude da corrente. Neste caso, a corrente e a tensão estão \emph{em fase}.

\subsubsection{Circuito Capacitivo Puro}

Num circuito capactitivo puro, com fonte de tensão alternada \(V(t) = V_0 \sin(\omega t)\), a 1ª Lei de Kirchhoff nos dá

\begin{equation}
    - V_0 \sin(\omega t) + \frac{q(t)}{C} = 0
\end{equation} onde \(q(t)\) é a carga no capacitor. Resolvendo para \(q(t)\), obtemos

\begin{equation}
    \boxed{q(t) = C V_0 \sin(\omega t)}
\end{equation} e a corrente no circuito é

\begin{equation}
    \boxed{i(t) = \frac{dq}{dt} = \omega C V_0 \cos(\omega t) = \frac{V_0}{X_C} \sin\left(\omega t + \frac{\pi}{2}\right) = I_0 \sin\left(\omega t + \frac{\pi}{2}\right)}
\end{equation} onde \(\displaystyle X_C \equiv 1/\omega C\) é a reatância capacitiva, de unidade ohm, e \(I_0 = \omega C V_0\) é a amplitude da corrente. A corrente está \emph{adiantada} de \(\pi/2\) em relação à tensão.

\subsubsection{Circuito Indutivo Puro}

Aplicando a lei das malhas para um circuito indutivo puro, com fonte de tensão alternada \(V(t) = V_0 \sin(\omega t)\), obtemos

\begin{equation}
    - V_0 \sin(\omega t) + L \frac{di(t)}{dt} = 0
\end{equation} onde \(i(t)\) é a corrente no indutor. Resolvendo para \(i(t)\), temos

\begin{equation}
    \boxed{i(t) = -\frac{V_0}{\omega L} \cos(\omega t) = \frac{V_0}{X_L} \sin\left(\omega t - \frac{\pi}{2}\right) = I_0 \sin\left(\omega t - \frac{\pi}{2}\right)}
\end{equation} onde \(\displaystyle X_L \equiv \omega L\) é a reatância indutiva, de unidade ohm, e \(I_0 = V_0/\omega L\) é a amplitude da corrente. A corrente está \emph{atrasada} de \(\pi/2\) em relação à tensão.

\subsubsection{Circuito RLC em Série}

Para um circuito RLC em série com uma fonte de tensão alternada \(V(t) = V_0 \sin(\omega t)\), temos a equação da malha

\begin{equation}
    - V_0 \sin(\omega t) + R i(t) + L \frac{di(t)}{dt} + \frac{1}{C} q(t) = 0 \implies R\frac{dq}{dt} + L \frac{d^2q}{dt^2} + \frac{1}{C} q = V_0 \sin(\omega t)
\end{equation} cuja solução permanente é

\begin{equation}
    \boxed{i(t) = I_0 \sin(\omega t - \varphi)}
\end{equation}

Temos que as tensões no resistor, indutor e capacitor são dadas por

\begin{align}
    &V_R(t) = R i(t) = R I_0 \sin(\omega t - \varphi) \\
    &V_L(t) = L \frac{di(t)}{dt} = \omega L I_0 \cos(\omega t - \varphi) = X_L I_0 \sin\left(\omega t - \varphi + \frac{\pi}{2}\right) \\
    &V_C(t) = \frac{1}{C} q(t) = -\frac{1}{\omega C} I_0 \cos(\omega t - \varphi) = -X_C I_0 \sin\left(\omega t - \varphi - \frac{\pi}{2}\right)
\end{align} com tensão máxima \(\displaystyle V_0 = \sqrt{V_R^2 + (V_L - V_C)^2}\) e \(\displaystyle I_0 = \frac{V_0}{Z}\), onde \(Z \equiv \sqrt{R^2 + (X_L - X_C)^2}\) é a impedância do circuito. O ângulo de fase \(\varphi\) entre a corrente e a tensão é dado por

\begin{equation}
    \tan(\varphi) = \frac{X_L - X_C}{R}
\end{equation} e a ressonância ocorre quando \(X_L = X_C\), ou seja, \(\omega = \omega_0 \equiv 1/\sqrt{LC}\), onde a impedância é mínima \(Z = R\) e a corrente máxima \(I_0 = V_0/R\).

\subsubsection{Potência Média e Valor Eficaz}

A potência instantânea no resistor \(R\) é

\begin{equation}
    p(t) = Ri^2(t)
\end{equation} enquanto que a potência média \(\overline{P}\) é dada por

\begin{equation}
    \overline{P} = \frac{1}{T} \int_0^T p(t) dt = R \frac{1}{T} \int_0^T i^2(t) dt = R\overline{i^2}
\end{equation} onde 

\begin{equation}
    \overline{i^2} \equiv \frac{1}{T} \int_0^T i^2(t) dt = \frac{I_0^2}{T} \int_0^T \sin^2(\omega t) dt = \frac{I_0^2}{2}
\end{equation}

Definimos o valor eficaz (RMS: \emph{root mean square}) da corrente como

\begin{equation}
    \boxed{I_{rms} \equiv \sqrt{\overline{i^2}} = \frac{I_0}{\sqrt{2}}}
\end{equation}

Similarmente, o valor eficaz da tensão é

\begin{equation}
    \boxed{V_{rms} \equiv \sqrt{\overline{v^2}} = \frac{V_0}{\sqrt{2}}}
\end{equation}

Assim, a potência média dissipada no resistor é

\begin{equation}
    \boxed{\overline{P} = R I_{rms}^2 = V_{rms} I_{rms} \cos(\varphi)}
\end{equation} onde \(\cos(\varphi)\) é o fator de potência, que mede a eficiência do circuito em converter energia elétrica em trabalho útil. Em circuitos puramente resistivos, \(\varphi = 0\) e \(\overline{P} = V_{rms} I_{rms}\).

\subsection{Equações de Maxwell}

\begin{itemize}
    \item [(a)] \textbf{Lei de Gauss para o campo elétrico}
        \begin{equation}
            \boxed{\oint_{S} \mathbf{E}\cdot d\mathbf{S} = \frac{Q_{int}}{\varepsilon_0}}
        \end{equation} onde \(Q_{int}\) é a carga interna à superfície \(S\).

    \item [(b)] \textbf{Lei de Gauss para o campo magnético}
        \begin{equation}
            \boxed{\oint_{S} \mathbf{B} \cdot d\mathbf{S} = 0}
        \end{equation} (não existem monopolos magnéticos).
    \item [(c)] \textbf{Lei de Faraday}
        \begin{equation}
            \boxed{\oint_{C} \mathbf{E} \cdot d\mathbf{l} = -\frac{d}{dt}\int_{S} \mathbf{B} \cdot d\mathbf{S}}
        \end{equation}
    \item [(d)] \textbf{Lei de Ampère-Maxwell}
        \begin{equation}
            \boxed{\oint_{C} \mathbf{B} \cdot d\mathbf{l} = \mu_0 \left( I_{cond} + I_{d} \right)}
        \end{equation} onde \(\displaystyle I_{d} = \varepsilon_0 \frac{d\Phi_E}{dt}\) é a corrente de deslocamento.
\end{itemize}



\end{document}